\section{Discussion}
\label{sec:discussion}

\subsection{Different kinds of bug, and why some evaded a linearity check}
The defects fixed in this paper are not structurally alike, and
that difference is itself instructive for anyone debugging a similar
commercial-AFE instrument. Defect~1 (Section~\ref{sec:rootcause}) is a
configuration-level mistake -- using a real, working AD5941 capability
(the HS loop) for the wrong purpose -- diagnosable by direct verification
against vendor reference code because the correct configuration already
existed, fully specified, in Analog Devices' own examples. Defect~2
(Section~\ref{sec:adcfix}) is a numerical-convention bug that a
check at the level of ``which loop is configured'' would not by
itself have caught; it was only found by inspecting raw ADC codes
alongside converted values and cross-checking the arithmetic against the
vendor's own formula. D-RCAL (Section~\ref{sec:otherdefects}) is
different again, and more dangerous than either: it is a constant
multiplicative scaling error, and Table~\ref{tab:rcalsweep} demonstrates
directly that such an error is invisible to a linearity check -- mean
$R^2$ moved by only $2\times10^{-6}$ across a $982\times$ span in absolute
error. A quantitative accuracy check against a traceable reference, not
just a straight-line fit, is what actually exposed it. D-STALE and
D-CCMD compound this: even after D-RCAL's declared value is corrected,
a runtime correction only reaches the measurement path if the command
that sets it writes the variable that path actually reads, and only takes
effect on the next sweep if the calibration that consumes it is not
frozen from power-up. All three had to be fixed together for a runtime
Rcal correction to work as documented. D-PGA, D-VZEROBUF, D-SIGN, and
D-WAKE were present despite this project's own OCP and legacy CV paths --
which use independently-correct conventions -- functioning normally
throughout, a direct illustration of the risk the duplicated,
per-method logic (Section~\ref{sec:oop}) carried: a single mistake, made
once, existed in three literal copies and had to be found and fixed three
times over. The shared base class introduced in this work
(Section~\ref{sec:oop}) is intended to prevent exactly that recurrence for
whatever the next defect turns out to be.

\subsection{What was a deliberate design choice, not a gap}
\label{sec:deliberate}
To be fair to this project's own architecture, several of
Table~\ref{tab:summary}'s differences from the vendor reference are
intentional and appropriate for this instrument's actual use case -- a
host-GUI-driven, live-configurable bench instrument -- not oversights.
Analog Devices' examples are meant to be recompiled per experiment,
driving the AD5941 through its on-chip sequencer and a wakeup timer so the
MCU and AFE can sleep between samples, which is a power-optimized design
appropriate for battery-powered, long-duration monitoring. This project
instead runs a plain MCU-side polling loop and exposes CA/SWV/DPV
parameters through the existing ASCII serial protocol, so a user (or the
host GUI) can change them live without reflashing -- a genuine usability
advantage for a benchtop instrument where power is not constrained and
simplicity of the MCU-side control flow matters more than microsecond-level
timing precision or battery life. Fixing D-SIGN's bias-computation
subtraction and collapsing the vendor's two-parameter
\texttt{Vzero}/\texttt{SensorBias} model into one runtime-settable value
plus one hardcoded constant (Section~\ref{sec:rootcause}) is the same kind
of reasonable simplification, made at the cost of not being able to reach
potentials that would require shifting \texttt{Vzero} itself without a
firmware constant change.

\subsection{Limitations}
\label{sec:limitations}
\begin{itemize}
\item \textbf{Sustained CA/SWV/DPV current does not track applied
  potential -- open, unfixed.} Section~\ref{sec:ca_dummy} shows the CA
  first-sample response is correctly signed and monotonic across five
  applied steps, and SWV/DPV (Section~\ref{sec:swv_dpv}) are smooth and
  bounded. Samples taken after the first $\approx$\SI{13}{\milli\second}
  of a sustained CA run, however, converge to a common current
  independent of the applied step -- a $\pm200$~mV step across the
  Randles cell's \SI{10.56}{\kilo\ohm} DC resistance should differ by
  roughly \SI{38}{\micro\ampere} and, in this failure mode, differs by
  less than \SI{3}{\nano\ampere}. All eight LPTIA feedback-resistor
  settings and both SW13 filter states (the \texttt{k} command,
  Section~\ref{sec:components}) were swept and changed nothing, which
  places the fault upstream of the sensing/conversion stage -- in
  establishing the LPPA/CE0 cell drive itself -- rather than in
  Eq.~\eqref{eq:current}'s conversion or in the TIA gain. The
  sequencer-driven CV path measures the same physical cell correctly
  (Section~\ref{sec:cv_dummy}), so this is specific to the direct-register
  CA/SWV/DPV classes, not a property of the AFE or the board. This paper
  reports this defect honestly, as found, rather than either silently
  omitting it or claiming these three methods are validated for sustained
  measurement; root-causing and fixing it is left to future work.
\item \textbf{A CV timeout requires a full reset, not just the D-WAIT2
  bound.} Bounding the CV wait (D-WAIT2) stops the instrument wedging
  indefinitely, but does not by itself restore the AFE to a fully working
  state: sweeps taken immediately after a reported \texttt{CV\_TIMEOUT}
  are reproducibly degraded ($R^2$ drops from $\approx0.999999$ to
  $\approx0.9999$ at unchanged accuracy) and recover only after a full
  reset. All data reported in Section~\ref{sec:validation} was collected
  on a freshly reset board for this reason.
\item \textbf{OOP maturity.} As detailed in Section~\ref{sec:oop}, this
  paper closes every gap the original audit raised -- encapsulating
  \texttt{C\_DataStorage}, bringing all six method classes (including
  cyclic voltammetry, via a new \texttt{C\_CV} wrapper) under one
  polymorphic \texttt{C\_ElectrochemicalMethod} hierarchy, and eliminating
  the triplicated configuration logic -- verified against this paper's own
  unchanged CV accuracy result after the refactor
  (Section~\ref{sec:validation}). This is not the same claim as a
  maximally-abstracted architecture: there remains exactly one inheritance
  level, no call site dispatches through a base-class pointer (each
  \texttt{Dispatch<Method>()} still names its concrete class), and
  \texttt{C\_DataStorage}'s new accessors are mechanical get/set pairs, not
  a richer domain interface. The paper's claim is scoped to what the audit
  actually measured -- encapsulation, inheritance/polymorphism, and
  duplication -- not to architectural maturity in general.
\item \textbf{EIS not re-validated.} This work's scope was the reported
  CA/SWV data-quality failure and the further defects it led to; EIS,
  which already used the HS loop appropriately for its AC excitation
  requirement, was not re-tested against the dummy cell in this work.
\item \textbf{RCAL's own tolerance is the next-largest unexplained
  residual.} The Randles-cell CV accuracy result
  (Section~\ref{sec:cv_dummy}) is repeatable to $0.0084\,\%$ but biased
  $+0.427\,\%$ high, a residual that exceeds the component tolerance
  stack ($\approx\pm0.11\,\%$) of the Randles cell's own parts. The
  dominant untested candidate is the fitted \SI{200}{\ohm} RCAL1's own
  tolerance, which enters the scale directly via the calibration this
  paper's fix now performs correctly and is not separable through that
  same measurement path. Characterizing RCAL1 against a traceable
  standard is the single highest-value next measurement.
\item \textbf{Sample timing precision.} This project's \texttt{delay()}-based
  sample interval is the requested period plus however long the
  measurement itself took, not the vendor's hardware-guaranteed exact
  output data rate; at low sample rates the discrepancy is proportionally
  larger. This is unrelated to, and not fixed by, the corrections
  reported in this paper.
\item \textbf{No on-chip sequencer.} The AD5941's dedicated fast
  single-shot transient sequence, meant for exactly the kind of fast
  capacitive decay a chronoamperometry step produces, is not used by this
  project's polling-loop architecture (Section~\ref{sec:deliberate}); a
  future revision adopting it would trade implementation complexity for
  finer time resolution and lower power, and may also be relevant to
  root-causing the sustained-current defect above.
\item \textbf{OCP.} Not independently validated on this baseline; OCP's
  calibration path uses a distinct code path from CA/SWV/DPV and was not
  in this paper's diagnostic scope.
\end{itemize}
